% Section 4.3, from lectures/L04/appendix/c_stack.md.

\section{The Stack, and How Deep It Goes}
\label{c4:sec:stack}

\subsection{One pointer, four users, no referee}
\label{c4:sec:users}
\code{SP} lives in \code{SPH:SPL} at data space \code{0x5E:0x5D}, it starts at \code{RAMEND}, and
it moves downwards. Four unrelated things use it and none of them coordinates with the others:

\begin{narrowtable}{@{}ll@{}}
  \thead{What} & \thead{Costs} \\
  \midrule
  \code{rcall} and \code{ret} & 2 bytes, for a return address \\
  \code{push} and \code{pop} & 1 byte each \\
  An interrupt & 2 bytes, pushed by the hardware before your handler starts \\
  Your handler's own saves & 1 byte each \\
\end{narrowtable}

\textbf{\code{SP} points at the next free byte, not the last used one.} So the number of bytes in
use is \code{RAMEND - SP}, and the lowest byte actually occupied is one \emph{above} where
\code{SP} is. That off-by-one is worth getting right once: with \code{SP} down at \code{0x08F0}, it
is the difference between ``my structure at \code{0x08F0} is safe'' and ``my structure at
\code{0x08F0} was overwritten''.

\subsection{Counting the worst case}
\label{c4:sec:worst}

\bookfigure{stack_depth}{Stack use accumulating in Chapter~3's program: nothing in the main loop,
which calls nothing, two bytes when an interrupt pushes the program counter, eleven after the
handler saves eight registers and \code{SREG}, thirteen and fifteen for the two nested calls inside
it.}{c4:fig:depth}

That is Chapter~3's program at its deepest, and every number in it is measured. The rule behind
them is simple enough to apply to any program:
\[
\begin{aligned}
  \text{bytes} = {} & 2 \times (\text{calls deep in the main program}) \\
    & {} + 2 + (\text{registers the handler saves}) \\
    & {} + 2 \times (\text{calls deep inside the handler})
\end{aligned}
\]
The $+\,2$ in the middle is the hardware's, and it is there even for a handler that is nothing but
\code{reti}. The handler's saves are one byte each, \code{SREG}'s holder included.

\textbf{The worst case is not the common case, and it is the only one worth computing.} The
interrupt does not arrive politely between main-loop iterations; it arrives while the main loop is
as deep as it ever gets, because that is the state the program spends its time in. Assume it lands
at the worst possible moment, because eventually it will.

\textbf{Everything is measurable except the worst moment.} You can watch the stack pointer and
record how low it went, and this book's harness does exactly that, but a measurement only tells
you about the interrupts that actually arrived and where they actually landed. The arithmetic tells
you about the one that has not happened yet. The cross-check in Exercise~\ref{c4:ex:crosscheck} is
about the two agreeing, and about what it would mean if the measurement came out \emph{lower}.

\subsection{What running out looks like}
\label{c4:sec:overflow}
Nothing.

There is no stack limit register, no guard page, no fault and no message. The stack grows down into
whatever is below it and keeps going. What you observe is one of these, days later:
\begin{itemize}
  \item a variable that changes on its own, because a \code{push} landed on it;
  \item a subroutine that returns somewhere it was never called from, because something wrote over
    a saved return address;
  \item a program that restarts for no reason, because that somewhere was unprogrammed flash
    (Exercise~\ref{c1:ex:program}, part d).
\end{itemize}

None of those points at the stack, which is what makes the arithmetic worth doing in advance rather
than the debugging worth doing afterwards.

\textbf{Recursion is the fast way there.} A handler that enables interrupts inside itself, or a
subroutine that calls itself, adds a frame per repetition with nothing counting them
(\secref{c3:sec:mustnot}). Neither appears in this book, and both appear in real AVR code.

\subsection{How much room is there really}
\label{c4:sec:room}
2048 bytes sounds like a lot until you write the subtraction down.

Chapter~3's program uses 15 bytes of stack at its worst and 17 bytes of structures, one LED and one
button, which leaves 2016 free. An array of six LEDs adds 42 bytes of structures, and a main loop
that walks it goes three calls deep, \code{led_array_all_on} into \code{led_on} into
\code{shift_bits}, so the figures move to 21 and 59, leaving 1968. Neither is close to trouble, and
that is the point of computing it: you want to know \emph{how much} headroom you have, so that when
a later change halves it you notice.

Do that subtraction whenever the program grows a call level or a handler, and write the answer
somewhere the next reader will find it. It is the kind of calculation that is done once, believed
for a year, and quietly falsified by the third feature nobody re-checked.

\subsection{The stack pointer is a register you may load}
\label{c4:sec:loadsp}
Everything above treats \code{SP} as something the hardware moves. It is also an ordinary pair of
I/O registers, and you may read it and write it like any other.

That is worth saying plainly, because \secref{c1:sec:program} writes \code{SP} once, at reset, and
then explains that the hardware had already done it, which leaves the impression that touching
\code{SP} is ceremonial. It is not. It is the only way to answer two questions this section has
been asking on paper.

\needspace{5\baselineskip}
\textbf{Reading it: how deep did I actually go?}

\begin{asm}
    in   r24, SPL
    in   r25, SPH          ; r25:r24 = SP, and RAMEND - SP is the depth
\end{asm}

\needspace{5\baselineskip}
\textbf{Writing it: high byte first.}

\begin{asm}
    out  SPH, r25
    out  SPL, r24
\end{asm}

The order matters and only in one direction. \code{SP} has no shadow register, and unlike
\code{OCR1A} both halves are directly writable, so the hazard is not a latch, it is the moment in
between. Write \code{SPL} first and there is one instruction during which the low half is new and
the high half is old: \code{SP} points somewhere neither value describes, and an interrupt arriving
in that instruction pushes two bytes there. Writing \code{SPH} first leaves the intermediate value
inside the \emph{new} region rather than in an arbitrary one. Interrupts being off makes the
question moot, and writing them in the safe order costs nothing, so the rule is worth following
rather than reasoning about.

\textbf{These are I/O addresses, and this section has been quoting the other kind.} \code{SPL} is
data space \code{0x5D} and I/O \code{0x3D}; \code{SPH} is \code{0x5E} and \code{0x3E}. \code{in}
and \code{out} take the I/O numbers, and \code{m328Pdef.inc} gives you the names, so
\code{in r24, SPL} is what you write (\secref{c1:sec:iowindow}). Reaching them with \code{lds} and
\code{sts} works too and costs a cycle more each.

\textbf{What this is for.} Nothing in this book needs it: the drivers you are writing never move
\code{SP}, and the depth figures in \secref{c4:sec:worst} come from the harness watching \code{SP}
from outside rather than from your program reading it. It is here because of what the four
instructions above become when you write them in the other order. Save every register on the stack
you are on, store \code{SP} somewhere, load \code{SP} from somewhere else, restore every register
from \emph{that} stack, and \code{ret}, and you have returned onto a stack you were not called on.
That is a context switch, it is about forty lines, and the two instruction pairs in this section are
the two that make it a switch rather than a subroutine.
